\documentclass[10pt,a4paper]{ctexart}
\usepackage[margin=1in]{geometry}
\usepackage{amsmath,amssymb}
\usepackage{booktabs}
\usepackage{enumitem}
\usepackage{xurl}
\usepackage[colorlinks=true,linkcolor=blue,urlcolor=blue]{hyperref}

\title{论文算法部分（中文版）}
\author{}
\date{\today}

\begin{document}
\maketitle

\section{算法}
\label{sec:algorithm_cn}

\subsection{概述}

我们提出 \textbf{TCDP}（\textit{Topology-Conditioned Delay Predictor}，拓扑条件延迟预测器）用于低资源跨工艺节点的时序弧级延迟预测。
对于每条待预测 timing arc，TCDP 融合两类互补输入：
1）描述工作条件与器件统计的紧凑数值特征；
2）由晶体管级电路图提取的弧感知异构图嵌入。
跨节点迁移由共享图感知主干实现，而跨域失配由“域校准残差 + 拓扑条件残差”联合修正。

\subsection{问题定义}

对第 $i$ 个样本，令 $c_i$ 为标准单元类型，
\begin{equation}
a_i = \left(p_i^{\mathrm{in}}, p_i^{\mathrm{out}}\right)
\end{equation}
为查询的输入到输出时序弧。
令
\begin{equation}
\omega_i = \left(\sigma_i,\kappa_i,\pi_i,s_i,l_i,v_i,T_i\right)
\end{equation}
分别表示时序方向/类型、归一化弧条件、极性、输入 slew、输出 load、电压与温度。
单个样本记为
\begin{equation}
\mathbf{u}_i = \left(c_i, a_i, \omega_i, y_i, d_i\right),
\end{equation}
其中 $y_i$ 为目标延迟，$d_i\in\{s,t\}$ 表示源域/目标域。

每个样本还关联一个数值描述向量
\begin{equation}
\mathbf{x}_i \in \mathbb{R}^{d_x},
\end{equation}
该向量由工作条件项、SPICE 提取的晶体管宽度统计、极性信息与轻量电气代理特征组成。
目标是学习 $\hat y_i = F(c_i,a_i,\omega_i,\mathbf{x}_i)$，并在跨节点场景下保持稳定泛化。

\subsection{寄生感知异构图编码器}

对每个标准单元，我们由晶体管级网表构建异构图
\begin{equation}
\mathcal{G}_{c}=\left(\mathcal{V}_{\mathrm{NET}}\cup\mathcal{V}_{\mathrm{PMOS}}\cup\mathcal{V}_{\mathrm{NMOS}},\mathcal{E}\right).
\end{equation}
图中节点和关系类型编码了门控、源漏连接、反向连接以及寄生相关交互。

令 $\mathbf{f}_v$ 为节点 $v$ 的原始特征，先做类型投影：
\begin{equation}
\mathbf{h}_v^{(0)} = W_{\phi(v)}\mathbf{f}_v.
\end{equation}
随后堆叠 $L$ 层异构注意力块：
\begin{equation}
\tilde{\mathbf{h}}_v^{(\ell+1)}=
\sum_{r\in\mathcal{R}(v)}
\mathrm{GAT}_r^{(\ell)}
\left(\left\{\mathbf{h}_u^{(\ell)}:u\in\mathcal{N}_r(v)\right\}\right),
\end{equation}
\begin{equation}
\mathbf{h}_v^{(\ell+1)}=
\mathrm{Dropout}\!\left(
\mathrm{ReLU}\!\left(
\mathrm{LN}\!\left(\tilde{\mathbf{h}}_v^{(\ell+1)}+\mathbf{h}_v^{(\ell)}\right)
\right)\right).
\end{equation}

读出阶段以查询弧 $a_i$ 为条件，聚焦其相关局部子结构：
\begin{equation}
\mathbf{z}_i = E_{\phi}\!\left(\mathcal{G}_{c_i},a_i\right).
\end{equation}

\subsection{拓扑条件延迟预测器延迟预测器}

将数值特征与图嵌入拼接：
\begin{equation}
\mathbf{q}_i = [\mathbf{x}_i \Vert \mathbf{z}_i], \qquad
\mathbf{h}_i = f_{\mathrm{bb}}(\mathbf{q}_i).
\end{equation}

最终预测分解为三项：
\begin{equation}
\hat y_i =
f_{\mathrm{sh}}(\mathbf{h}_i)
+ f^{\mathrm{cal}}_{d_i}(\mathbf{h}_i)
+ f^{\mathrm{topo}}_{d_i}\!\left([\mathbf{h}_i \Vert \mathbf{e}_{\tau_i}]\right),
\end{equation}
其中 $\tau_i$ 是由同逻辑族（合并驱动强度变体）得到的拓扑组，$\mathbf{e}_{\tau_i}$ 为其可学习嵌入。
第二项用于域相关校准，第三项用于拓扑族相关残差修正。

\subsection{跨节点学习目标}

在 epoch $e$，令 $B_t,B_s$ 分别为目标域和源域小批样本。采用 MSE：
\begin{equation}
\mathcal{L}_t = \frac{1}{|B_t|}\sum_{i\in B_t}(\hat y_i-y_i)^2,\qquad
\mathcal{L}_s = \frac{1}{|B_s|}\sum_{i\in B_s}(\hat y_i-y_i)^2.
\end{equation}
训练目标写为
\begin{equation}
\mathcal{L}_{\mathrm{batch}}(e)=\mathcal{L}_t+\alpha(e)\mathcal{L}_s,
\end{equation}
其中源域权重按退火策略更新：
\begin{equation}
\alpha(e)=
\begin{cases}
1, & e<e_a,\\
1-(1-\alpha_f)\dfrac{e-e_a}{e_b-e_a}, & e_a\le e<e_b,\\
\alpha_f, & e\ge e_b.
\end{cases}
\end{equation}
该策略在训练早期保留源域可迁移监督，后期逐步强调目标域拟合。

\subsection{消融映射（与实验节一致）}

为确保“组件变化”与“性能变化”一一对应，算法与消融映射如下：
\begin{equation}
\texttt{TopologyOff}: \ f^{\mathrm{topo}}_{d}(\cdot)\equiv 0,
\end{equation}
\begin{equation}
\texttt{DualMean}: \ \mathbf{z}_i=\frac{\mathbf{z}_i^{\mathrm{global}}+\mathbf{z}_i^{\mathrm{arc}}}{2},
\end{equation}
\texttt{ParasiticReplace} 与 \texttt{EnrichOff(auto)} 对应图编码输入阶段的 NET 特征构造变化。
该映射保证了 one-factor 消融的可解释性与可复现性。

总体而言，TCDP 不依赖复杂分布对齐模块，而是通过“共享图感知主干 + 域校准残差 + 拓扑条件残差”完成跨节点适配，兼顾结构归纳偏置与工程实现轻量性。

\end{document}

